\documentclass[11pt,a4paper]{article}

% ---- Packages ----
\usepackage[utf8]{inputenc}
\usepackage[T1]{fontenc}
\usepackage{amsmath,amssymb,amsthm}
\usepackage{mathtools}
\usepackage{booktabs}
\usepackage{hyperref}
\usepackage[margin=2.5cm]{geometry}
\usepackage{enumitem}
\usepackage{graphicx}
\usepackage{multirow}

% ---- Theorem environments ----
\newtheorem{theorem}{Theorem}
\newtheorem{proposition}[theorem]{Proposition}
\newtheorem{conjecture}[theorem]{Conjecture}
\newtheorem{definition}[theorem]{Definition}
\newtheorem{corollary}[theorem]{Corollary}
\newtheorem{lemma}[theorem]{Lemma}
\theoremstyle{remark}
\newtheorem{remark}[theorem]{Remark}

% ---- Shortcuts ----
\newcommand{\pb}[2]{\{#1,\,#2\}}
\newcommand{\R}{\mathbb{R}}
\newcommand{\N}{\mathbb{N}}
\newcommand{\Z}{\mathbb{Z}}
\newcommand{\cA}{\mathcal{A}}
\newcommand{\cL}{\mathcal{L}}
\newcommand{\cF}{\mathcal{F}}
\newcommand{\cT}{\mathcal{T}}
\newcommand{\cJ}{\mathcal{J}}
\newcommand{\SG}{S_3}
\DeclareMathOperator{\GKdim}{GKdim}
\DeclareMathOperator{\Span}{span}
\DeclareMathOperator{\rank}{rank}

\title{$S_3$-equivariant jet filtration of the\\
three-body Poisson algebra:\\
tier decomposition, integer scaling exponents,\\
and syzygy structure}

\author{[Author names]}

\date{\today \quad \textit{(Draft)}}

\begin{document}
\maketitle

\begin{abstract}
In a companion paper we showed that the Poisson algebra
$\cA$ generated by the three pairwise Hamiltonians of the
planar three-body problem has dimension sequence
$d(0)=3$, $d(1)=6$, $d(2)=17$, $d(3)=116$,
with 156 candidate generators at level~3 of which 116
are linearly independent.  Here we reveal the internal
organisation of this algebra.  The singular value (SV)
spectrum at each configuration exhibits a four-tier
staircase
\[
  52 + 44 + 16 + 4 = 116
\]
with inter-tier gaps spanning $10^{2}$--$10^{4}$ orders
of magnitude.  We prove that the tier sizes are determined
by the $\SG$ isotypic decomposition of the algebra:
the 156 generators carry exactly 24~copies of the
trivial representation~$A$, 28~copies of the sign
representation~$A'$, and 52~copies of the standard
representation~$E$, and the number $n_E = 52$ equals
the first tier size exactly.  The Clebsch--Gordan rules
of $\SG$ predict the multiplicities at every bracket level.

We identify a jet filtration underlying the tier
structure: fitting singular values as $\sigma \sim
\varepsilon^\alpha$ across five sampling radii yields
integer-quantized scaling exponents
$\alpha = 0, 1, 2, 3$ for tiers 1--4 respectively,
with the noise floor flat at machine epsilon.
The tiers thus correspond to successive orders of the
Taylor expansion of generator functions in configuration
space.  The gap at the rank boundary grows as
$\sim\!\varepsilon^2$, from $5.5\times$ at
$\varepsilon=10^{-4}$ to $67{,}472\times$ at
$\varepsilon=2\times 10^{-3}$.

The 40 null directions in the SV spectrum are not
Dirac constraints but syzygies: all 156 generators are
individually nonzero as phase-space functions, and the
null space decomposes into 8~universal identities
(machine-epsilon everywhere) and 32~soft syzygies
that break at the collinear submanifold (raising the
rank from 116 to~125).  The tier decomposition and the
syzygy structure are orthogonal.  All results are
potential-independent: the $1/r$ and $1/r^2$ systems
share the same tier sizes, scaling exponents, and null
counts.
\end{abstract}


% ====================================================================
\section{Introduction}
% ====================================================================

In a companion paper~\cite{paper1} we studied the Poisson
algebra~$\cA$ generated by the three pairwise interaction
Hamiltonians $H_{12}$, $H_{13}$, $H_{23}$ of the planar
three-body problem.  The cumulative dimension through
bracket level~$n$,
\[
  d(0) = 3,\quad d(1) = 6,\quad d(2) = 17,\quad d(3) = 116,
\]
grows super-exponentially, implying infinite
Gelfand--Kirillov dimension.  The sequence is invariant
under changes of mass ratios and identical for the
Newtonian ($1/r$) and Calogero--Moser ($1/r^2$) potentials,
establishing it as an algebraic invariant of the singularity
class.

That analysis characterised the algebra by its
\emph{dimension}.  The present paper asks a finer question:
\emph{how is the 116-dimensional space internally
organised?}  We find three interlocking structures:

\begin{enumerate}[nosep]
  \item A \textbf{four-tier decomposition}
    $52 + 44 + 16 + 4 = 116$, visible as shelves in the
    singular value spectrum separated by gaps of
    $10^{2}$--$10^{4}$ (Section~\ref{sec:tiers}).
  \item An \textbf{$\SG$ isotypic decomposition} that
    explains the tier sizes exactly via Clebsch--Gordan
    rules (Section~\ref{sec:s3}).
  \item A \textbf{jet filtration} in which the tiers
    correspond to orders $0,1,2,3$ of the local Taylor
    expansion, with integer-quantized scaling exponents
    (Section~\ref{sec:jets}).
\end{enumerate}

We also resolve the status of the 40 null directions in
the SV spectrum: they are syzygies among iterated Poisson
brackets, not Dirac constraints
(Section~\ref{sec:syzygies}).  All structures are
universal across singular potentials and configurations.


% ====================================================================
\section{Setup and notation}
\label{sec:setup}
% ====================================================================

We follow the conventions of~\cite{paper1}.  Three point
masses $m_1, m_2, m_3$ in the plane have phase space
$\R^{12}$ with canonical coordinates
$(x_i, y_i, p_{x_i}, p_{y_i})$, $i = 1,2,3$.  The pairwise
Hamiltonians are
\[
  H_{ij} = \frac{p_{x_i}^2 + p_{y_i}^2}{2m_i}
          + \frac{p_{x_j}^2 + p_{y_j}^2}{2m_j}
          + V(r_{ij}),
\]
where $V(r) = -1/r$ (Newtonian) or $V(r) = -1/r^2$
(Calogero--Moser).  Using the polynomial representation
$u_{ij} = 1/r_{ij}$, all Poisson brackets are computed
in exact rational arithmetic (levels~0--3) or numerically
(level~4).

\subsection{Generator census}

The algebra~$\cA$ is generated by $\{H_{12}, H_{13},
H_{23}\}$ under iterated Poisson brackets.  At each
level~$n$, the new generators are all brackets
$\pb{f}{g}$ with $f$ from the frontier at level~$n-1$
and $g$ from the accumulated set through level~$n-1$,
modulo antisymmetry.  The candidate counts are:
\begin{center}
\begin{tabular}{@{}ccrr@{}}
\toprule
Level & Construction & Candidates & New generators \\
\midrule
0 & $H_{12}, H_{13}, H_{23}$ & 3 & 3 \\
1 & $\pb{\cA_0}{\cA_0}$ & 3 & 3 \\
2 & $\pb{\cA_1}{\cA_0} + \Lambda^2\cA_1$ & 12 & 11 \\
3 & $\pb{\cA_2}{\cA_0} + \pb{\cA_2}{\cA_1}
    + \Lambda^2\cA_2$ & 138 & 99 \\
\midrule
\multicolumn{2}{c}{Total} & 156 & 116 \\
\bottomrule
\end{tabular}
\end{center}
Of the 156 candidate generators, 116 are linearly
independent at generic phase-space points and 40 satisfy
algebraic relations.  The number 156 is the total
candidate count (all non-antisymmetric bracket pairs);
it is not the number of nonzero expressions.

\subsection{Numerical evaluation}

To probe the algebra at a specific configuration~$q$ in
shape space, we sample $N$ phase-space points uniformly
in an $\varepsilon$-ball around~$q$ and evaluate all 156
generators to form the evaluation matrix
$M \in \R^{N \times 156}$.  After column
normalisation, the singular value decomposition
$M = U\Sigma V^T$ yields the sorted singular values
$\sigma_1 \ge \sigma_2 \ge \cdots \ge \sigma_{156}$.
The effective rank is determined by the largest gap in the
$\sigma$-spectrum.

We parametrise the reduced configuration space (shape
sphere for equal masses) by the mass ratio
$\mu = m_3 / m_1$ and the angle~$\phi$ measuring
departure from collinear alignment.  A $100 \times 100$
grid over $(\mu, \phi) \in [0.2, 3.0] \times [3^\circ,
177^\circ]$ provides 10{,}000 configurations, each
evaluated at five sampling radii
$\varepsilon \in \{10^{-4}, 2\times 10^{-4},
5\times 10^{-4}, 10^{-3}, 2\times 10^{-3}\}$.


% ====================================================================
\section{The tier decomposition}
\label{sec:tiers}
% ====================================================================

\subsection{Observation: a universal staircase}

At every non-degenerate configuration, the normalised SV
spectrum $\sigma_k / \sigma_1$ displays a characteristic
staircase pattern with four shelves separated by abrupt
drops.

\begin{proposition}[Tier structure]
\label{prop:tiers}
For both the $1/r$ and $1/r^2$ potentials, at every
configuration with $\phi > 10^\circ$, the $156$-component
SV spectrum evaluated at $\varepsilon = 10^{-4}$ exhibits
four tiers:
\begin{center}
\begin{tabular}{@{}clcc@{}}
\toprule
Tier & SV indices & Size & $\sigma/\sigma_1$ range \\
\midrule
1 & $0$--$51$ & 52 & $1$ to $10^{-1}$ \\
2 & $52$--$95$ & 44 & $10^{-3}$ to $10^{-4}$ \\
3 & $96$--$111$ & 16 & $10^{-7}$ to $10^{-8}$ \\
4 & $112$--$115$ & 4 & $10^{-11}$ to $10^{-13}$ \\
\midrule
Noise & $116$--$155$ & 40 & $< 10^{-15}$ \\
\bottomrule
\end{tabular}
\end{center}
The inter-tier gaps span $10^{2}$--$10^{4}$ in the SV
ratio, making the decomposition unambiguous.
\end{proposition}

Table~\ref{tab:tier-configs} shows the tier boundaries at
representative configurations for both potentials.  The
first three boundaries (positions 52, 96, 112) are
remarkably stable: they agree between $1/r$ and $1/r^2$
at every tested configuration.

\begin{table}[ht]
\centering
\small
\begin{tabular}{@{}llcc@{}}
\toprule
Configuration & Potential & Tier positions & Subspace sizes \\
\midrule
Lagrange ($\mu{=}1, \phi{=}60^\circ$)
  & $1/r$   & 52, 96, 110, 115 & 52, 44, 14, 5 \\
  & $1/r^2$ & 52, 94, 110, 115 & 52, 42, 16, 5 \\[3pt]
Isosceles ($\mu{=}1, \phi{=}90^\circ$)
  & $1/r$   & 52, 96, 112, 115 & 52, 44, 16, 3 \\
  & $1/r^2$ & 52, 96, 112, 114 & 52, 44, 16, 2 \\[3pt]
Large-$\mu$ Lagrange
  & $1/r$   & 52, 96, 112, 116 & 52, 44, 16, 4 \\
  & $1/r^2$ & 52, 96, 112, 116 & 52, 44, 16, 4 \\[3pt]
Collinear ($\phi \approx 0^\circ$)
  & $1/r$   & 151 & 151 \\
  & $1/r^2$ & 151 & 151 \\
\bottomrule
\end{tabular}
\caption{Tier boundaries at special configurations
  ($\varepsilon = 10^{-4}$).  The ``ideal'' decomposition
  $52 + 44 + 16 + 4$ emerges cleanly at large-$\mu$
  Lagrange; minor variations ($\pm 2$ generators) occur
  at exact symmetry points.  The collinear configuration
  is anomalous: all 151 candidate generators become
  significant.}
\label{tab:tier-configs}
\end{table}

\subsection{Statistical significance}

Across the full $100 \times 100$ atlas (10{,}000 grid
points, both potentials, 48{,}318 total tier detections),
the preferred tier positions cluster sharply at indices
52, 96, 116, and~129.  The Kolmogorov--Smirnov test
against a uniform distribution gives $D = 0.667$,
$p < 10^{-4}$.

The subspace-size distribution has concentration index
(one minus normalised entropy) $c = 0.254$.  In
10{,}000 Monte Carlo trials with uniformly random sizes
in $[1, 156]$, the maximum concentration achieved was
$c \approx 0.18$.  The observed quantization of tier
sizes is statistically extreme ($p < 10^{-4}$).

The most frequent subspace sizes (Table~\ref{tab:sizes})
coincide with the tier sizes and their complements to~116.

\begin{table}[ht]
\centering
\begin{tabular}{@{}crr@{}}
\toprule
Size & Count & Fraction \\
\midrule
116 & 13{,}758 & 28.5\% \\
4   & 5{,}538  & 11.5\% \\
112 & 2{,}625  & 5.4\% \\
52  & 1{,}202  & 2.5\% \\
44  & --- & --- \\
16  & 1{,}374  & 2.8\% \\
\bottomrule
\end{tabular}
\caption{Most frequent subspace sizes across 48{,}318
  tier detections.  The size~116 dominates because at many
  configurations only the main cliff is detected.}
\label{tab:sizes}
\end{table}


% ====================================================================
\section{$S_3$ representation theory}
\label{sec:s3}
% ====================================================================

The symmetric group $\SG$ acts on the three-body
configuration space by permuting the labels
$(1,2,3)$.  For equal masses this is a genuine symmetry
of the Hamiltonian; for unequal masses it is a symmetry
of the algebraic structure (which monomials appear in the
generators), even though the numerical coefficients
differ~\cite{paper1}.

\subsection{Character table and Clebsch--Gordan rules}

$\SG$ has three irreducible representations:
\begin{center}
\begin{tabular}{@{}lccc@{}}
\toprule
Irrep & $\dim$ & Character on $(12)$ & Character on $(123)$ \\
\midrule
$A$ (trivial) & 1 & 1 & 1 \\
$A'$ (sign) & 1 & $-1$ & 1 \\
$E$ (standard) & 2 & 0 & $-1$ \\
\bottomrule
\end{tabular}
\end{center}

The Clebsch--Gordan decomposition rules are:
\begin{align*}
  A \otimes A &= A, &
  A \otimes A' &= A', &
  A \otimes E &= E, \\
  A' \otimes A' &= A, &
  A' \otimes E &= E, &
  E \otimes E &= A \oplus A' \oplus E.
\end{align*}
The exterior square relevant for self-brackets is
$\Lambda^2(E) = A'$.

\subsection{Isotypic decomposition of the algebra}

The three pairwise Hamiltonians $\{H_{12}, H_{13},
H_{23}\}$ form the edge representation of $\SG$, which
decomposes as $A \oplus E$.  At level~1, the three tidal
operators $K_i = \pb{H_{jk}}{H_{j\ell}}$ (one per
vertex) also form $A \oplus E$.  Higher levels are
generated by the Clebsch--Gordan rules.

\begin{theorem}[Isotypic decomposition]
\label{thm:isotypic}
The 156 candidate generators of~$\cA$ through bracket
level~3 decompose under $\SG$ as
\[
  24\,A \;\oplus\; 28\,A' \;\oplus\; 52\,E
  \qquad (\text{total dimension } 24 + 28 + 104 = 156).
\]
At each bracket level, the decomposition is:
\begin{center}
\begin{tabular}{@{}crrrrr@{}}
\toprule
Level & $n_A$ & $n_{A'}$ & $n_E$ & $\dim$
  & $E$-fraction \\
\midrule
0 & 1 & 0 & 1 & 3 & 66.7\% \\
1 & 1 & 0 & 1 & 3 & 66.7\% \\
2 & 2 & 2 & 4 & 12 & 66.7\% \\
3 & 20 & 26 & 46 & 138 & 66.7\% \\
\midrule
Total & 24 & 28 & 52 & 156 & 66.7\% \\
\bottomrule
\end{tabular}
\end{center}
The $E$-fraction is exactly $2/3$ at every bracket level.
\end{theorem}

\begin{proof}
The seed $\cA_0 = A \oplus E$ has $E$-fraction $2/3$.
The key observation is that the Clebsch--Gordan
decomposition of an $\SG$-module with $E$-fraction $2/3$
preserves this fraction.  At each level~$n$, new
generators arise from two sources:
\begin{enumerate}[nosep]
  \item \emph{Cross-brackets} $\pb{\cA_{n-1}}{\cA_k}$
    for $k < n-1$: the tensor product of two modules
    with $E$-fraction $2/3$ again has $E$-fraction $2/3$
    (verified by direct computation using the CG rules).
  \item \emph{Self-brackets} $\Lambda^2(\cA_{n-1})$:
    the exterior square $\Lambda^2(nA \oplus mA' \oplus
    \ell E)$ has $E$-fraction $2/3$ when the input has
    $E$-fraction $2/3$ and $n + m = \ell$ (which holds at
    every level by induction, since $n_A + n_{A'} = n_E$
    for the seed and the CG rules preserve this).
\end{enumerate}
We verified the decomposition computationally:
\begin{itemize}[nosep]
  \item Level~2:
    $\pb{\cA_1}{\cA_0} = (A{+}E) \otimes (A{+}E)
    = 2A + A' + 3E$ (dim~9),
    $\Lambda^2(\cA_1) = \Lambda^2(A{+}E) = A' + E$
    (dim~3), total $2A + 2A' + 4E = 12$.
    \textbf{Matches observed count.}
  \item Level~3:
    $\pb{\cA_2}{\cA_0} = 6A + 6A' + 12E$ (dim~36),
    $\pb{\cA_2}{\cA_1} = 6A + 6A' + 12E$ (dim~36),
    $\Lambda^2(\cA_2) = 8A + 14A' + 22E$ (dim~66),
    total $20A + 26A' + 46E = 138$.
    \textbf{Matches observed count.}
    \qedhere
\end{itemize}
\end{proof}

\subsection{The number 52}

\begin{corollary}
\label{cor:52}
The first tier size equals the number of $E$-type
irreducible copies:
\[
  |\text{Tier 1}| = n_E = 52.
\]
\end{corollary}

This identity is exact, not approximate.  The 52
$E$-doublets contribute $2 \times 52 = 104$ generators
to the algebra.  At a generic (non-symmetric)
configuration, the $\SG$ symmetry is broken: each
$E$-doublet splits into two singlets with distinct
singular values.  The more dynamically significant
component of each doublet enters Tier~1, producing
exactly 52 entries.

\begin{remark}[Doublet signatures]
At the near-Lagrange configuration
($\mu = 0.992$, $\phi = 60.2^\circ$), the $\SG$ symmetry
is approximately restored.  We detect near-degenerate
doublets (SV ratio~$< 1.05$) in the spectrum:
\begin{center}
\begin{tabular}{@{}lrrr@{}}
\toprule
Tier & Doublets & Singlets & Total \\
\midrule
1 & 18 & 16 & 52 \\
2 & 11 & 22 & 44 \\
3 & 2 & 12 & 16 \\
4 & 0 & 4 & 4 \\
\bottomrule
\end{tabular}
\end{center}
The high doublet count in Tier~1 confirms the
$E$-doublet origin.  Full degeneracy would require exact
$\SG$ symmetry; the splitting reflects the grid's
approximation to the Lagrange point.
\end{remark}

\subsection{Tier structure as broken isotypic decomposition}

At the $\SG$-symmetric Lagrange point, the $E$-type
generators carry the most dynamical information because
they sense the \emph{relative geometry} of the three
bodies---they transform nontrivially under permutations.
The $A$-type generators measure \emph{total} (symmetric)
quantities, and $A'$-type generators measure
\emph{antisymmetric} combinations.

The tier structure reflects this hierarchy:
\begin{itemize}[nosep]
  \item \textbf{Tier~1 (52)}: One component from each of
    the 52 $E$-doublets---the dynamically dominant
    direction.
  \item \textbf{Tier~2 (44)}: The partner component from
    44 of the 52 doublets.  The 8 missing partners have
    fallen below the Tier~2 boundary (severe symmetry
    breaking displaces one partner).
  \item \textbf{Tiers~3 and~4 ($16 + 4 = 20$)}: A subset
    of the $n_A + n_{A'} = 24 + 28 = 52$ generators in
    one-dimensional representations.
  \item \textbf{Noise (40)}: The remaining
    $52 - 20 = 32$ generators from $A$/$A'$ irreps and
    8~displaced $E$-partners, collectively forming the
    syzygy space (Section~\ref{sec:syzygies}).
\end{itemize}

\begin{remark}
The decomposition $n_A + n_{A'} = 52 = n_E$ is not a
coincidence: it follows from the CG rules applied to the
seed $A \oplus E$ and the fact that the $E$-fraction is
locked at $2/3$ at every level.  Since each $E$ contributes
dimension~2, the $A/A'$ count $n_A + n_{A'}$ equals the
$E$ count $n_E$.
\end{remark}


% ====================================================================
\section{Jet filtration and integer scaling}
\label{sec:jets}
% ====================================================================

The tier decomposition has a natural interpretation in
terms of jet spaces.  Define the $k$-jet of a smooth
function $f$ at a configuration~$q$ as its Taylor
polynomial through order~$k$:
\[
  j^k_q f = \sum_{|\alpha| \le k}
    \frac{1}{\alpha!}
    \left(\frac{\partial^{|\alpha|} f}
    {\partial z^\alpha}\right)\!\bigg|_q
    (z - q)^\alpha,
\]
where $z$ denotes phase-space coordinates.  The
evaluation matrix $M$ constructed by sampling in an
$\varepsilon$-ball captures the Taylor expansion of each
generator up to terms of order~$O(\varepsilon^K)$ for
some effective truncation order~$K$.  The singular values
of $M$ should then scale as $\sigma_k \sim
\varepsilon^{\alpha_k}$ where $\alpha_k$ is the lowest
Taylor order at which the $k$-th independent direction
contributes new information.

\subsection{Scaling exponents}

\begin{theorem}[Integer scaling]
\label{thm:scaling}
At the Lagrange configuration, fitting $\log_{10}
\sigma_k$ against $\log_{10} \varepsilon$ across five
sampling radii gives scaling exponents
$\alpha = \alpha_k$ that are integer-quantized:
\begin{center}
\begin{tabular}{@{}lcccc@{}}
\toprule
Component & Count & $\bar\alpha \pm \delta\alpha$
  & Nearest integer & Interpretation \\
\midrule
Tier 1 & 52 & $-0.007 \pm 0.014$ & 0
  & zeroth-order observables \\
Tier 2 & 44 & $\phantom{-}1.002 \pm 0.025$ & 1
  & first-order variation \\
Tier 3 & 16 & $\phantom{-}2.005 \pm 0.165$ & 2
  & second-order variation \\
Tier 4 & 4 & $\phantom{-}2.824 \pm 0.210$ & 3
  & third-order variation \\
Noise  & 40 & flat at $\sim\!10^{-15}$
  & $\infty$ & identically zero \\
\bottomrule
\end{tabular}
\end{center}
The noise floor shows no $\varepsilon$-dependence because
the corresponding linear combinations are exactly zero.
\end{theorem}

\begin{proof}[Proof (numerical)]
For each SV index~$k$ and each of the five sampling radii
$\varepsilon_j$, the singular value $\sigma_k(\varepsilon_j)$
is recorded at the Lagrange configuration
($\mu = 1, \phi = 60^\circ$).  Least-squares regression
of $\log_{10}\sigma_k$ against $\log_{10}\varepsilon_j$
yields the slope~$\alpha_k$ and coefficient of
determination~$R^2$.  For indices $0 \le k \le 115$
(the 116 significant directions), $R^2 > 0.99$ in all
cases, confirming the power-law relationship.  The
mean exponents by tier are as stated above.

For the noise floor ($k \ge 116$), the singular values
are at machine epsilon ($\sim\!10^{-16}$) regardless of
$\varepsilon$, confirming that these directions are
identically zero.
\end{proof}

\begin{remark}[Jet filtration]
Define the $k$-th jet subspace $\cJ^k_q \subset \cA_3$ as
the span of generators whose local Taylor expansion around
$q$ begins at order~$\le k$.  Theorem~\ref{thm:scaling}
shows that
\[
  \cJ^0_q \subset \cJ^1_q \subset \cJ^2_q
  \subset \cJ^3_q = \cA_3|_q
\]
with
\[
  \dim \cJ^0_q = 52, \quad
  \dim \cJ^1_q = 96, \quad
  \dim \cJ^2_q = 112, \quad
  \dim \cJ^3_q = 116.
\]
The successive quotients have dimensions $52, 44, 16, 4$,
matching the tier sizes.  This jet filtration refines the
bracket-level filtration $\cA_0 \subset \cA_1 \subset
\cA_2 \subset \cA_3$ (which has dimensions $3, 6, 17,
116$) and is orthogonal to it: generators from the same
bracket level can appear in different jets.
\end{remark}

\subsection{Gap growth}

The gap ratio at the rank boundary (SV index~116) grows
with $\varepsilon$:

\begin{center}
\begin{tabular}{@{}cc@{}}
\toprule
$\varepsilon$ & $\sigma_{115}/\sigma_{116}$ \\
\midrule
$10^{-4}$ & $5.5$ \\
$2 \times 10^{-4}$ & $51$ \\
$5 \times 10^{-4}$ & $950$ \\
$10^{-3}$ & $8{,}100$ \\
$2 \times 10^{-3}$ & $67{,}472$ \\
\bottomrule
\end{tabular}
\end{center}

The growth is approximately $\propto \varepsilon^2$,
consistent with Tier~4 generators scaling as
$\varepsilon^3$ while the noise floor remains constant.
At $\varepsilon \ge 10^{-3}$ the rank determination is
definitive (gap $> 10^3$).

\subsection{Relationship between jet and bracket filtrations}

The bracket-level filtration and the jet filtration are
distinct.  The bracket filtration has dimensions
$3, 6, 17, 116$ and is purely algebraic (determined by
bracket depth).  The jet filtration has dimensions
$52, 96, 112, 116$ and depends on the \emph{local}
Taylor expansion around a configuration.

Levels~0 and~1 (6~generators total) contribute only to
Tier~1, as expected: the Hamiltonians and their first
brackets are zeroth-order observables.  Level~2 generators
(11~new) appear in Tiers~1 and~2.  Level~3 generators
(99~new) populate all four tiers.  The jet filtration thus
provides finer information than the bracket depth.


% ====================================================================
\section{Syzygy structure of the null space}
\label{sec:syzygies}
% ====================================================================

The 40 directions in the SV spectrum below the rank-116
cliff require careful interpretation.  If they
corresponded to generators that vanish identically on
phase space, they would be Dirac-type constraints and the
tier structure might reflect a first-class/second-class
distinction rather than an isotypic decomposition.

\subsection{No generator vanishes individually}

\begin{proposition}
\label{prop:no-constraints}
All 156 candidate generators of $\cA_3$ are individually
nonzero as functions on $\R^{12}$.  At every tested
configuration (Lagrange, isosceles, collinear, scalene)
and every sampling radius
$\varepsilon \in \{10^{-4}, 10^{-3}, 10^{-2}\}$, the
RMS value of each generator exceeds $10^{-14}$.
\end{proposition}

\begin{proof}[Proof (computational)]
The symbolic Poisson algebra was built through level~3,
producing 156~generators as symbolic expressions.  Each
was evaluated independently at 500 random phase-space
points in an $\varepsilon$-ball around each of four
special configurations.  The column RMS
$(\sum_i g(z_i)^2 / N)^{1/2}$ exceeded $10^{-14}$ for
all generators at all configurations and all~$\varepsilon$.
\end{proof}

The 40 null directions in the SVD are therefore
\emph{syzygies}: linear combinations
$\sum_k c_k\, g_k = 0$ of generators that cancel,
not individual generators that vanish.

\subsection{Deep and soft syzygies}

The 40 syzygies decompose into two populations:

\begin{definition}
A syzygy is \emph{deep} if the corresponding singular
value is at machine epsilon ($< 10^{-14}$) at every
configuration and every $\varepsilon$.  A syzygy is
\emph{soft} if it holds at generic configurations but
breaks at special submanifolds.
\end{definition}

\begin{proposition}
\label{prop:syzygy-decomp}
The 40-dimensional null space of $\cA_3$ decomposes as
\[
  \underbrace{8\text{ deep syzygies}}_{\text{universal
    identities}} \;\oplus\;
  \underbrace{32\text{ soft syzygies}}_{\text{Jacobi
    consequences}}.
\]
\begin{enumerate}[nosep]
  \item The 8 deep syzygies (SV indices 148--155) satisfy
    $\sigma < 2.3 \times 10^{-14}$ across all
    10{,}000 configurations and all 5 sampling radii.
    Indices 151--155 are at machine epsilon
    ($< 10^{-15}$) everywhere.

  \item The 32 soft syzygies (SV indices 116--147) hold
    at generic configurations but break near the
    collinear submanifold ($\phi \approx 6^\circ$): 8 of
    these ``wake up,'' increasing the effective rank
    from 116 to~125 at $(\mu = 1.020,
    \phi = 5.7^\circ)$.
\end{enumerate}
\end{proposition}

\begin{remark}
The soft syzygies are consequences of the Jacobi
identity $\pb{A}{\pb{B}{C}} + \text{cyc} = 0$.  At
generic configurations, the three cyclic terms cancel.
At collinear configurations, where certain brackets
degenerate, the cancellation fails and formerly dependent
directions become independent.
\end{remark}

\subsection{Null-space structure}

Analysis of the right singular vectors $v_k$ for
$k \ge 116$ reveals:

\begin{enumerate}[nosep]
  \item \textbf{Level concentration.}  96.4\% of the
    null-space weight is in level-3 generators,
    $\sim\!3.4\%$ in level~2, and 0\% in levels~0--1.
    No seed or first-bracket generator participates in
    any syzygy.

  \item \textbf{Cross-level mixing.}  30--39 of the
    40 null vectors mix generators from different bracket
    levels.  These are \emph{cross-level syzygies}:
    algebraic identities relating generators at different
    bracket depths.

  \item \textbf{Potential independence.}  Both $1/r$ and
    $1/r^2$ give exactly 40 null directions at every
    configuration.  The syzygies are a property of the
    bracket structure and the three-body interaction
    graph, not the potential.
\end{enumerate}

\subsection{Orthogonality of tiers and syzygies}

The tier decomposition and the syzygy structure are
orthogonal:
\[
\underbrace{156\text{ generators}}_{\text{all nonzero}}
= \underbrace{116\text{ independent}}_{\text{organised
  into 4 tiers by } \SG}
\;\oplus\; \underbrace{40\text{ syzygies}}_{\text{algebraic
  identities}}.
\]
The tier boundaries reflect \emph{dynamical significance}
mediated by $\SG$ representation theory: which generators
carry the most geometric information at a given
configuration.  The syzygy boundary reflects
\emph{algebraic redundancy}: which generators are
expressible as linear combinations of others.  These are
independent classifications.

\begin{proposition}
\label{prop:rank-stability}
The modal rank is 116 at all thresholds $\ge 10^{-10}$.
The rank is stable:
\begin{center}
\begin{tabular}{@{}cccc@{}}
\toprule
Threshold & Min rank & Modal rank & Max rank \\
\midrule
$10^{-8}$ & 108 & 112 & 116 \\
$10^{-10}$ & 114 & 116 & 125 \\
$10^{-12}$ & 116 & 116 & 138 \\
$10^{-13}$ & 116 & 116 & 148 \\
\bottomrule
\end{tabular}
\end{center}
Configurations with rank $> 116$ cluster near the
collinear submanifold where soft syzygies break.
\end{proposition}


% ====================================================================
\section{Discussion}
\label{sec:discussion}
% ====================================================================

\subsection{Intrinsic quantization}

The tier sizes 52, 44, 16, 4 do not correspond to any
known quantization scheme.  We tested explicitly:

\begin{itemize}[nosep]
  \item \textbf{Loop quantum gravity} (area spectrum
    $\sqrt{j(j{+}1)}$ for half-integer $j$):
    mean matching error 0.153, Monte Carlo $p = 0.77$
    (not significant).
  \item \textbf{Bekenstein} (equally-spaced integer
    areas): mean matching error 0.093,
    $p = 0.18$ (not significant).
  \item \textbf{$\mathrm{SU}(2)$ representation dimensions}
    ($2j{+}1$): no correspondence.
\end{itemize}

The quantization is intrinsic: it arises from the
$\SG$ isotypic decomposition of the Poisson algebra,
not from any external quantum-gravitational structure.

\subsection{Physical interpretation of E-dominance}

The $E$-type generators are dynamically dominant because
they sense the \emph{relative geometry} of the three
bodies.  In the standard representation $E$, $\SG$ acts by
permuting the vertices of an equilateral triangle; the
$E$-components of a function measure how the function
value changes when the particle labels are permuted.  At
a generic (scalene) configuration where all three bodies
are distinguishable, these are the most informative
observables.

The $A$-type generators measure \emph{total} quantities
invariant under permutation, and $A'$-type generators
measure fully antisymmetric combinations.  Both are less
informative at generic configurations because they do not
distinguish which body is which.

\subsection{Open questions}

\begin{enumerate}
  \item \textbf{Level~4 tier structure.}  Does the jet
    filtration extend to level~4?  The lower bound
    $d(4) \ge 4{,}501$ from~\cite{paper1} implies a
    vastly larger algebra.  Whether the four-tier
    pattern persists, or whether new jet orders appear,
    is unknown.

  \item \textbf{$N = 4$ analogue.}  The four-body
    dimension sequence $[6, 14, 62]$ (through level~2)
    has been computed~\cite{paper1}.  The relevant
    symmetry group is~$S_4$, with five irreps.  Whether
    an analogous tier decomposition exists---and whether
    its sizes are predicted by $S_4$
    Clebsch--Gordan rules---is a natural question.

  \item \textbf{Jet-order growth.}  The exponents
    $\alpha = 0, 1, 2, 3$ match the bracket levels
    $0, 1, 2, 3$ in number but not in dimension
    (the bracket filtration has dimensions $3, 6, 17, 116$
    while the jet filtration has dimensions $52, 96, 112,
    116$).  Is there a representation-theoretic formula
    predicting $\dim \cJ^k_q$ in terms of the $\SG$
    content at each bracket level?

  \item \textbf{Syzygy resolution.}  The 8 deep syzygies
    likely correspond to consequences of translational
    symmetry (total momentum conservation).  Can they be
    identified explicitly as expressions proportional
    to $P_x$ and $P_y$ brackets at various levels?

  \item \textbf{Collinear anomaly.}  Why do 8 soft
    syzygies break specifically at near-collinear
    configurations?  The collinear arrangement is an
    $S_2$-symmetric stratum (one pair of bodies defines
    a distinguished axis); the syzygy-breaking may
    reflect the reduction $\SG \to S_2$ and the
    branching rules for the sign representation.
\end{enumerate}


% ====================================================================
\section*{Acknowledgments}
% ====================================================================

The computational pipeline---symbolic Poisson bracket
engine, SVD analysis, atlas scanning, and statistical
tests---was developed with the assistance of Claude
(Opus~4.6)~\cite{claude-2026}, a large language model by
Anthropic.  All mathematical results were verified by
independent symbolic and numerical computation.


% ====================================================================
\begin{thebibliography}{99}

\bibitem{paper1}
[Author names],
\emph{Super-exponential growth of the Poisson algebra
generated by pairwise Hamiltonians of the planar
three-body problem},
Preprint, 2026.

\bibitem{claude-2026}
Anthropic,
\emph{Claude (Opus~4.6)} [Large language model], 2026.
\url{https://claude.ai/}

\bibitem{serre-1977}
J.-P.~Serre,
\emph{Linear Representations of Finite Groups},
Springer, 1977.

\bibitem{fulton-harris-1991}
W.~Fulton and J.~Harris,
\emph{Representation Theory: A First Course},
Springer, 1991.

\bibitem{kollar-2007}
J.~Koll\'{a}r,
\emph{Lectures on Resolution of Singularities},
Ann.\ of Math.\ Stud., vol.\ 166, Princeton Univ.\
Press, 2007.

\bibitem{moser-1975}
J.~Moser,
\emph{Three integrable Hamiltonian systems connected with
isospectral deformations},
Adv.\ Math.\ \textbf{16} (1975), 197--220.

\bibitem{calogero-1971}
F.~Calogero,
\emph{Solution of the one-dimensional $N$-body problems
with quadratic and/or inversely quadratic pair potentials},
J.~Math.\ Phys.\ \textbf{12} (1971), 419--436.

\bibitem{dirac-1950}
P.~A.~M.~Dirac,
\emph{Generalized Hamiltonian dynamics},
Canad.\ J.~Math.\ \textbf{2} (1950), 129--148.

\bibitem{henneaux-teitelboim-1992}
M.~Henneaux and C.~Teitelboim,
\emph{Quantization of Gauge Systems},
Princeton Univ.\ Press, 1992.

\end{thebibliography}

\end{document}
